%% Pitts 1995 Remark 5.6.2: the parameter-free equality bi-rule.
\calculussheet{sheet:contextual-parameter-free-equality}{Pitts's parameter-free equality rule}{These are the two directions of the bi-rule in \cite[Remark~5.6.2, p.~56]{Pitts1995Draft}, denoted locally by \(\mathrm{Eq0}\). The variable contexts displayed are the complete contexts: there is no suppressed parameter prefix. The variables \(x,x'\) are distinct and have sort \(\sigma\). The binary formula abstraction \(\psi\) has body \(\psi(x,x')\) well formed in exactly \((x:\sigma,x':\sigma)\). The empty hypothesis list is denoted \(\varnothing\). Bound variables are freshened and substitution avoids capture. Theorem~\ref{thm:p10-contextual-equality-reduction} proves equivalence to the parameterized equality rule over the stated common basis with implication and universal quantification.}
\hbrule{rule:contextual-parameter-free-equality-down}{Eq0-down}{\varnothing\vdash\psi(x,x)\,[x:\sigma]}{(x=_\sigma x')\vdash\psi(x,x')\,[x:\sigma,x':\sigma]}{\(\psi\) has no free parameters other than its two argument variables.}
\hbrule{rule:contextual-parameter-free-equality-up}{Eq0-up}{(x=_\sigma x')\vdash\psi(x,x')\,[x:\sigma,x':\sigma]}{\varnothing\vdash\psi(x,x)\,[x:\sigma]}{\(\psi\) has no free parameters other than its two argument variables.}
